\itemhead{p.21 练习 · 第1题（教材p.21）}{XJ-B1-C01-PR-P021-Y0793}
\label{XJ-B1-C01-PR-P021-Y0793}
{\footnotesize 湘教版教材p.21 / PDF第28页 / p.21 练习 / 第1题}
\field{原题}{指出下列命题中使用了什么量词以及量词的作用范围，并把量词用相应的数学符号取代：
\subq{1}对任意整数 $n$，$2n+1$ 是奇数；
\subq{2}对某个有理数 $x$，有 $4x=\dfrac12$；
\subq{3}线段 $AB$ 上有一点 $M$ 满足比例式 $\dfrac{AM}{MB}=\dfrac{MB}{AB}$。}
\field{参考答案}{（1）全称，$\forall n\in\mathbb Z,2n+1$ 为奇数；（2）存在，$\exists x\in\mathbb Q,4x=\frac12$；（3）存在，$\exists M\in\overline{AB}$，满足 $MB\ne0$ 且 $AM/MB=MB/AB$。}
{\small \field{关联知识}{L03 全称、存在量词及真假的判定（教材帮p.38、39 / PDF47、48）}
\field{方法与易错点}{先确认量词、作用对象及论域，再采用任意性证明或实例反例。 易错：用几个满足的数证明全称命题；把“存在”理解为“唯一存在”。}
\field{答案依据}{答案册相关题组：PDF第6页（册内页5）；按当前题库原题逐项重算，不按旧版题号机械对套。}}
